% =============================================================================
% Kapitel 4: Methodik der vorliegenden Arbeit
% =============================================================================
\chapter{Methodik}
\label{ch:methodik}

\section{Proben und Messdaten}

Die Arbeit behandelt zwei Probensysteme, deren EXAFS-Daten im Rahmen
der Vorlesung bereitgestellt wurden:

\begin{description}[style=nextline]
  \item[Probe~1: \TODO{CoNiCu}-Mittelentropielegierung (MEA)]
    Eine äquimolare ternäre Legierung der 3d-Übergangsmetalle Cobalt,
    Nickel und Kupfer.
    Das Absorptionsspektrum wurde in einem einzigen Energiescan von
    ca.\ \SIrange{7500}{10000}{\electronvolt} aufgenommen und deckt
    damit die K-Kanten aller drei Elemente ab
    (Co: \SI{7709}{\electronvolt}, Ni: \SI{8333}{\electronvolt},
    Cu: \SI{8979}{\electronvolt}).
  \item[Probe~2: \TODO{CoPt}-Nanopartikel]
    Bimetallische Nanopartikel aus Cobalt und Platin.
    \TODO{Hier werden nach Erhalt der Daten die Messbedingungen,
    Kanten und Energiebereiche ergänzt.}
\end{description}

\section{Datenvorverarbeitung}
\label{sec:datenvorverarbeitung}

\subsection{$\chi(k)$-Extraktion}

Die Umwandlung der Rohdaten $\mu(E)$ in $\chi(k)$ wurde mit der
Python-Bibliothek \textsc{larch} \cite{Newville2013} durchgeführt.
Die wesentlichen Schritte sind:

\begin{enumerate}
  \item \textbf{Kantenenergie $E_0$:} Bestimmung als Wendepunkt
        der Absorptionskante.
  \item \textbf{Pre-Edge-Subtraktion:} Anpassung einer linearen
        Funktion im Bereich unterhalb der Kante und Abzug.
  \item \textbf{Normierung:} Anpassung eines Polynoms im
        Post-Edge-Bereich und Normierung auf den Kantensprung
        $\Delta\mu_0 = 1$.
  \item \textbf{Hintergrundentfernung:} Spline-Fit
        (\texttt{autobk}-Algorithmus mit $R_{\text{bkg}} = \SI{1.0}{\Angstrom}$)
        zur Bestimmung und Subtraktion des atomaren Hintergrunds
        $\mu_0(E)$.
\end{enumerate}

\subsection{Besonderheit bei Multi-Edge-Messungen}

Da bei der MEA-Probe alle drei Kanten in einem einzigen Scan vorliegen,
wurde für jede Kante ein separates Energiefenster definiert, um eine
Kontamination der Normierung durch die nächstfolgende Kante zu
vermeiden:
\begin{itemize}[nosep]
  \item Co-K: bis \SI{8233}{\electronvolt} (ca.\ \SI{100}{\electronvolt}
        vor der Ni-Kante),
  \item Ni-K: \SIrange{8100}{8879}{\electronvolt},
  \item Cu-K: ab \SI{8700}{\electronvolt}.
\end{itemize}
Die korrekte Normierung ist kritisch, da ein fehlerhafter Kantensprung
$\Delta\mu_0$ die gesamte Amplitude von $\chi(k)$ skaliert und damit
die effektiven Koordinationszahlen verfälscht.

\section{Startstrukturen}

Für die \evax{}-Analyse wurden dreidimensionale Superzellen mit
jeweils 256~Atomen erzeugt.
Die chemischen Elemente wurden stöchiometrisch korrekt und zufällig
auf die Gitterplätze verteilt.

Drei verschiedene Kristallstrukturen dienen als Startmodelle:
\begin{itemize}[nosep]
  \item \textbf{FCC} (kubisch-flächenzentriert): $4\times4\times4$
        konventionelle Einheitszellen, Gitterparameter
        $a = \SI{3.54}{\Angstrom}$, Superzellkante
        $L = \SI{14.16}{\Angstrom}$.
  \item \textbf{BCC} (kubisch-raumzentriert): $4\times4\times8$
        Einheitszellen, $a = \SI{2.87}{\Angstrom}$.
  \item \textbf{HCP} (hexagonal dichteste Packung):
        $8\times4\times4$ Einheitszellen,
        $a = \SI{2.51}{\Angstrom}$, $c = \SI{4.07}{\Angstrom}$.
\end{itemize}
Durch den Vergleich der finalen Residuen lässt sich bestimmen, welche
Struktur die experimentellen Daten am besten beschreibt.

\section{EvAX-Konfiguration}

\subsection{Fit-Parameter}

Die optimalen Fit-Parameter wurden durch eine systematische
Parameterstudie (vgl.\ Abschnitt~\ref{sec:param-scan}) ermittelt.
Tabelle~\ref{tab:fit-params} fasst die finalen Einstellungen zusammen.

\begin{table}[htbp]
  \centering
  \caption{Finale Fit-Parameter der \evax{}-Analyse.}
  \label{tab:fit-params}
  \small
  \begin{tabular}{ll}
    \toprule
    \textbf{Parameter} & \textbf{Wert} \\
    \midrule
    $k$-Bereich        & \SIrange{3}{11.5}{\per\Angstrom} \\
    $k$-Gewichtung     & $k^2$ \\
    Vergleichsraum     & Wavelet (\texttt{Space = w}) \\
    $N_{\text{legs}}$  & 4 (Mehrfachstreuung) \\
    $R_{\text{max,FEFF}}$ & \SI{6}{\Angstrom} \\
    Schrittweite       & \SI{0.005}{\Angstrom} \\
    Max.\ Verschiebung & \SI{0.4}{\Angstrom} \\
    Anzahl States      & 16 \\
    Screening-Iterationen & 2000 \\
    Froze-in           & 5000 \\
    \bottomrule
  \end{tabular}
\end{table}

\subsection{Automatische $S_0^2$- und $\Delta E_0$-Bestimmung}

Die Parameter $S_0^2$ und $\Delta E_0$ wurden über die
\texttt{-autoES}-Option von \evax{} automatisch bestimmt.
\evax{} verwendet dafür eine Vorstufe, in der diese Werte für jede
Kante separat optimiert werden, bevor die eigentliche
Strukturoptimierung beginnt.
Physikalisch sinnvolle Werte liegen bei $S_0^2 \approx 0{,}6$--$0{,}9$
und $|\Delta E_0| \lesssim \SI{15}{\electronvolt}$.

\section{Systematische Parameterstudie}
\label{sec:param-scan}

Um die Sensitivität der Ergebnisse gegenüber der Wahl der
Fit-Parameter zu untersuchen, wurde eine automatisierte
Parameterstudie durchgeführt.
In einer ersten Phase wurden die Fit-Parameter $k_{\min}$,
$k_{\max}$, $k^n$, der Vergleichsraum und $N_{\text{legs}}$
systematisch variiert (insgesamt 72~Kombinationen).
In einer zweiten Phase wurden die Optimierungsparameter
(Schrittweite, maximale Verschiebung, $R_{\max}^{\text{FEFF}}$)
mit den besten Fit-Parametern aus Phase~1 kombiniert
(11~Kombinationen).

Alle Läufe wurden parallelisiert (sechs gleichzeitige EvAX-Instanzen
in isolierten Arbeitsverzeichnissen) und mit reduzierten
Iterationszahlen (50~Screening-Iterationen, \texttt{Froze\_in}~=~50)
durchgeführt, um die Rechenzeit auf ein praktikables Maß zu begrenzen.
Die Ergebnisse dienen der Identifikation optimaler
Parametereinstellungen und der Abschätzung der Parametersensitivität,
nicht der finalen Strukturbestimmung.

\section{Strukturvergleich}
\label{sec:struct-compare}

Für die Bestimmung der zugrunde liegenden Kristallstruktur werden
Produktionsläufe mit den optimalen Parametern für jede der drei
Ausgangsstrukturen (FCC, BCC, HCP) durchgeführt.
Die Runs verwenden volle Konvergenzkriterien
(2000~Screening-Iterationen, \texttt{Froze\_in}~=~5000) und werden
anhand folgender Kriterien verglichen:

\begin{enumerate}[nosep]
  \item Gesamtresidual und kantenaufgelöste Einzelresiduen,
  \item Übereinstimmung der berechneten und experimentellen
        $\chi(k)$- und FT-Spektren,
  \item Physikalische Plausibilität der resultierenden $S_0^2$-Werte,
  \item Konsistenz der RDF-Peaks mit bekannten Bindungslängen.
\end{enumerate}

\section{Auswertung der Strukturergebnisse}

Aus den finalen \evax{}-Strukturen werden folgende Größen berechnet:

\begin{itemize}[nosep]
  \item \textbf{Partielle RDF} $g_{ij}(R)$ für alle Elementpaare,
  \item \textbf{Koordinationszahlen} $N_{ij}$ durch Integration des
        ersten Peaks,
  \item \textbf{Mittlere Bindungslängen} $\langle R_{ij}\rangle$ aus
        der Peakposition,
  \item \textbf{Warren-Cowley-Parameter} $\alpha_{ij}$
        (Gl.~\eqref{eq:warren_cowley}).
\end{itemize}

Die Unsicherheiten werden aus der Streuung der Werte über die
16~parallelen States abgeschätzt.
